\begin{itemize}
\item \textbf{$N=10$: cross-validation picks the simplest model (K=1 in 93\% of repetitions).} Each training split has only nine points, and component 4 has on average only two samples in the whole dataset. A two-dimensional full-covariance Gaussian has 5 free parameters, so every extra component multiplies the estimation variance. Extra components typically lock onto one or two training points with tiny covariances, and the held-out point then falls far from them. The median validation log-likelihood falls from $-6.0$ at $K=1$ to $-81$ at $K=4$ and below $-1000$ at $K\ge8$, more than two orders of magnitude worse than the true density's $-3.90$. A single broad Gaussian covering all three clusters is the best predictive model this little data can support. Ten samples carry essentially no information about the true order.
\item \textbf{$N=100$: the two overlapping components are merged (K=3 in 92\%).} With 100 samples the three \emph{visible} clusters are well estimated, which explains the jump in validation log-likelihood from $K=2$ to $K=3$. Components 1 and 2, however, have means only $1.71$ apart, comparable to the sum of their average eigenvalues ($0.90+0.85=1.75$). Their combined density is a single unimodal, slightly elongated lump, and roughly 55 points in it are not enough to estimate the 6 extra parameters a split requires. In 92 of 100 repetitions the 4-component model scored lower on held-out data than the 3-component model; the median difference was $-0.065$ nats per sample. The bias of merging the pair costs less than the variance of separating it.
\item \textbf{$N=1000$: the true order dominates (K=4 in 83\%).} With roughly 550 points in the overlapping pair, the 4-component model's better fit outweighs its extra variance. $K=4$ beat $K=3$ on held-out likelihood in 97\% of repetitions, by a median of $0.013$ nats per sample, and the median CV curve peaks at $K=4$ ($-3.918$), within $0.02$ of the true-density value $-3.899$. The remaining selections mostly lean toward \emph{over}fitting ($K=5$: 13\%, $K=6$: 3\%). The validation curve beyond $K=4$ is very flat (median drop of only $0.004$ from $K=4$ to $5$), because a superfluous component can sit almost redundantly inside an existing one at little cost. With that little separation, sampling noise sometimes tips the choice upward. Underfitting ($K=3$) occurs only once.
\item \textbf{Overall pattern.} The selected order rises with the amount of data (mean $\hat K$: 1.08, 3.08, 4.18). Cross-validated log-likelihood estimates the model with the best \emph{predictive} density for the data at hand, not the number of generating components. With limited data it deliberately prefers simpler models, and it recovers the true structure only once the data can resolve it. The overlapping pair is exactly where the true order is hardest to identify, and it accounts for nearly all the $K=3$ choices at $N=100$.
\item \textbf{Statistical precision.} Each rate is a proportion out of $R=100$ independent repetitions, with standard error $\sqrt{r(1-r)/100}\le0.05$. For example, the $N=1000$ rate for $K=4$ is $0.83\pm0.04$. The ranking of the most frequent order at every $N$ is therefore well beyond Monte Carlo noise. Rates of 0\% mean ``not observed in 100 trials'', not ``impossible''.
\item \textbf{Sensitivity to settings.} The covariance floor ($10^{-3}I$) and the 3 EM restarts are part of the fitted model family. A larger floor would soften the large-$K$ collapse at $N=10$; the conclusions at $N=100$ and $N=1000$ come from well-populated components, whose covariance eigenvalues are hundreds of times larger than the floor, so they are not driven by it.
\end{itemize}
